% Preâmbulo compartilhado dos artefatos LaTeX da disciplina de Aprendizagem
% de Máquina (graduação e pós-graduação).
%
% Regras editoriais: sem emojis, sem travessões, português do Brasil com
% acentuação correta. Arquivos em UTF-8 sem BOM.
%
% Conjunto de pacotes deliberadamente enxuto para garantir compilação com
% pdflatex em diferentes instalações (TeX Live, MiKTeX), sem tikz nem
% microtype.

\usepackage[utf8]{inputenc}
\usepackage[T1]{fontenc}
\usepackage[brazilian]{babel}
\usepackage{geometry}
\geometry{a4paper,margin=2.5cm}
\usepackage{graphicx}
\usepackage{booktabs}
\usepackage{xcolor}
\usepackage{enumitem}
\usepackage{caption}
\usepackage[hidelinks]{hyperref}

% Cores institucionais aproximadas (CIn-UFPE).
\definecolor{cinazul}{RGB}{0, 51, 102}
\definecolor{cinverde}{RGB}{0, 102, 51}

% Macro para título de etapa do projeto: Etapa 1 (teórica) e Etapa 2
% (experimental), com entrada correspondente no sumário.
\newcommand{\etapa}[2]{%
    \section*{Etapa #1: #2}%
    \addcontentsline{toc}{section}{Etapa #1: #2}%
}

% Macro para destacar tarefa pendente do grupo em uma caixa visível.
\newcommand{\tarefagrupo}[1]{%
    \par\medskip
    \noindent\fbox{%
        \begin{minipage}{\dimexpr\textwidth-2\fboxsep-2\fboxrule\relax}
        \textbf{Tarefa do grupo:} #1
        \end{minipage}%
    }%
    \par\medskip
}

% Macro para citar arquivos do repositório em fonte monoespaçada.
\newcommand{\arquivo}[1]{\texttt{#1}}

\hypersetup{
    pdfauthor={Prof. Leandro Maciel Almeida},
    pdfkeywords={aprendizagem de máquina, classificação tabular, projeto final}
}
